\subsection*{(1) Compute $VaR_{0.95}(L)$}

The cumulative distribution function (CDF) of $L$ is given by:
\begin{equation*}
F_L(x) = 
\begin{cases} 
0 & \text{if } x < 0 \\
0.95 & \text{if } 0 \le x < 100 \\
1 & \text{if } x \ge 100 
\end{cases}
\end{equation*}

By definition, $VaR_{\alpha}(L) = \inf\{x \in \mathbb{R} : F_L(x) \ge \alpha\}$. 
For $\alpha = 0.95$, the smallest value $x$ such that $F_L(x) \ge 0.95$ is $x = 0$. 
Therefore,
\begin{equation*}
VaR_{0.95}(L) = 0
\end{equation*}


\subsection*{(2) Compute $\mathbb{E}[L|L\ge VaR_{0.95}(L)]$}

Since we found $VaR_{0.95}(L) = 0$, we need to compute $\mathbb{E}[L | L \ge 0]$.

Because the only possible values for $L$ are 0 and 100, the condition $L \ge 0$ is always satisfied (it holds with probability 1). Thus, the conditional expectation is simply the unconditional expectation:
\begin{align*}
\mathbb{E}[L | L \ge 0] &= \mathbb{E}[L] \\
&= 0 \times \mathbb{P}(L = 0) + 100 \times \mathbb{P}(L = 100) \\
&= 0 \times 0.95 + 100 \times 0.05 \\
&= 5
\end{align*}


\subsection*{(3) Compute $ES_{0.95}(L)$}

The Expected Shortfall is defined via the integral:
\begin{equation*}
ES_{\alpha}(L) = \frac{1}{1-\alpha}\int_{\alpha}^1 VaR_u(L) du
\end{equation*}

For $u \in (0.95, 1]$, we are looking for $\inf\{x : F_L(x) \ge u\}$. Since $u > 0.95$, we must have $F_L(x) = 1$, which first occurs at $x = 100$. Thus, $VaR_u(L) = 100$ for all $u \in (0.95, 1]$.

Now we compute the integral for $\alpha = 0.95$:
\begin{align*}
ES_{0.95}(L) &= \frac{1}{1 - 0.95} \int_{0.95}^1 100 \, du \\
&= \frac{1}{0.05} \left[ 100u \right]_{0.95}^1 \\
&= 20 \times 100(1 - 0.95) \\
&= 20 \times 5 \\
&= 100
\end{align*}


\subsection*{(4) Explain why the conditional expectation definition is problematic here}

In this example, the conditional expectation yields $5$, whereas the true Expected Shortfall yields $100$. This shows a severe discrepancy.

The formula $ES_{\alpha}(L) = \mathbb{E}[L | L \ge VaR_{\alpha}(L)]$ is only guaranteed to hold when the cumulative distribution function $F_L$ is continuous at $VaR_{\alpha}(L)$. Here, the distribution is discrete, and there is a massive probability mass exactly at $VaR_{0.95}(L) = 0$, since $\mathbb{P}(L=0) = 0.95$.

When we condition on the event $L \ge VaR_{0.95}(L)$, we are inadvertently including the entire probability mass of the non-tail event $L=0$. This "dilutes" the expectation with a large number of zero-loss scenarios, fundamentally misrepresenting the average of the strictly worst 5\% of cases. The true integral definition of ES correctly isolates the 5\% tail, which in this distribution consists entirely of the $L=100$ outcome.